% apendices/mapa-particulas.tex
\chapter{🗺️ Mapa de Partículas}

\begin{tcolorbox}[vocabbox, title={📌 Partículas esenciales}]
\begin{tabularx}{\linewidth}{llX}
\toprule
\textbf{Part.} & \textbf{Función principal} & \textbf{Ejemplo} \\
\midrule
は (wa) & Tema del discurso & \ruby{私}{わたし}\textbf{は}\ruby{学生}{がくせい}です。 \\
が (ga) & Sujeto gramatical / énfasis & \ruby{猫}{ねこ}\textbf{が}います。 \\
を (wo) & Objeto directo & \ruby{本}{ほん}\textbf{を}\ruby{読}{よ}む。 \\
に (ni) & Destino / existencia / tiempo & \ruby{学校}{がっこう}\textbf{に}\ruby{行}{い}く。 \\
で (de) & Lugar de acción / medio & \ruby{図書館}{としょかん}\textbf{で}\ruby{勉強}{べんきょう}する。 \\
へ (e) & Dirección (más literario que に) & \ruby{日本}{にほん}\textbf{へ}。 \\
の (no) & Posesión / modificación & \ruby{私}{わたし}\textbf{の}\ruby{本}{ほん}。 \\
と (to) & Y (enumeración) / con / cita & \ruby{田中}{たなか}さん\textbf{と}。 \\
も (mo) & También / ni... tampoco & \ruby{私}{わたし}\textbf{も}。 \\
か (ka) & Pregunta & \ruby{学生}{がくせい}です\textbf{か}？ \\
よ (yo) & Afirmación / información nueva & そうです\textbf{よ}。 \\
ね (ne) & Confirmación / acuerdo & いいです\textbf{ね}。 \\
から (kara) & Desde / porque & \ruby{東京}{とうきょう}\textbf{から}。 \\
まで (made) & Hasta & \ruby{駅}{えき}\textbf{まで}。 \\
より (yori) & Más que (comparación) & A\textbf{より}B。 \\
\bottomrule
\end{tabularx}
\end{tcolorbox}

\begin{tcolorbox}[grammarbox, title={⚠️ は vs が — La distinción clave}]
\begin{tabular}{lll}
\toprule
\textbf{Uso} & \textbf{は} & \textbf{が} \\
\midrule
Introducir tema & ✓ & ✗ \\
Nueva información & ✗ & ✓ \\
Respuesta a ¿quién? & ✗ & ✓ \\
Contraste & ✓ & ✗ \\
Con 好き・嫌い・できる & ✗ & ✓ \\
\bottomrule
\end{tabular}\\[6pt]
\ej{\ruby{猫}{ねこ}\textbf{は}\ruby{魚}{さかな}が\ruby{好}{す}きです。}{Los gatos (tema) les gustan los pescados (objeto).}
\ej{\ruby{猫}{ねこ}\textbf{が}います。}{Hay un gato. (nueva información, énfasis)}
\end{tcolorbox}
